\documentclass[a4paper,landscape]{article}
\usepackage[margin=0mm,nohead,nofoot,nomarginpar]{geometry}
\usepackage{tikz}

% CONFIGURABLES
\def\linespacing{7}     % Spacing between ruled lines (mm)
\def\gridspacing{7}     % Square grid size (mm)
\def\dotspacing{5}      % Spacing between dots (mm)
\def\dotsize{0.15}      % Dot radius (mm)
% Ink darkness (0-100) and stroke width. A monochrome laser printer has no
% true gray - light tints get halftoned and thin/small features can print
% patchy. Solid (or near-solid) ink at a thin width prints crisper than a
% light gray fill; tune width for lightness instead of shade. Test-print
% and adjust to taste.
\def\shadeval{100}      % Ink darkness percentage
\def\strokewidth{0.15}  % Line width for ruled/grid/dot marks (mm)
\def\pagemargin{10}     % Outer side margin (mm)
\def\gutter{3}           % Extra margin on the spine side, for binding (mm)
\def\vtopmargin{10}      % Top margin (mm)
\def\vbottommargin{15}   % Bottom margin (mm) - leaves room for the page number
\def\sewingmargin{12}   % Sewing marks stay this far from the top/bottom (mm)
\def\startpage{1}       % Starting page number
\def\numsheets{20}      % A4 sheets to generate (each sheet = 2 A5 pages,
                        % single-sided) - 20 sheets = 40 A5 pages

% Calculate centering offset
\def\startx{0mm}
\edef\secondx{\the\dimexpr\startx+148.5mm\relax}

% Create A5 page content
% #1=x-offset, #2=y-offset, #3=type, #4=page number
% #5=spine side: 0 = spine on the right (left-hand page of the spread),
%                1 = spine on the left (right-hand page of the spread)
\newcommand{\makepage}[5]{%
    \begin{scope}[shift={(#1,#2)}]
        % Page border (optional)
        \draw[black!10] (0,0) rectangle (148.5mm,210mm);

        % Margins: the spine side gets an extra \gutter for binding
        \ifnum#5=0
            \edef\marginleft{\pagemargin}
            \edef\marginright{\the\numexpr\pagemargin+\gutter\relax}
            \edef\dotmarginleft{\dotspacing}
            \edef\dotmarginright{\the\numexpr\dotspacing+\gutter\relax}
        \else
            \edef\marginleft{\the\numexpr\pagemargin+\gutter\relax}
            \edef\marginright{\pagemargin}
            \edef\dotmarginleft{\the\numexpr\dotspacing+\gutter\relax}
            \edef\dotmarginright{\dotspacing}
        \fi

        % Content based on type
        \ifnum#3=1% Ruled
            \pgfmathtruncatemacro{\nlines}{floor((210-\vtopmargin-\vbottommargin)/\linespacing)}
            \foreach \i in {0,...,\nlines} {
                \draw[black!\shadeval,line width=\strokewidth mm]
                    (\marginleft mm,{210mm-\vtopmargin mm-\i*\linespacing mm}) --
                    ({148.5mm-\marginright mm},{210mm-\vtopmargin mm-\i*\linespacing mm});
            }
        \fi
        \ifnum#3=2% Grid
            % The first corner given is guaranteed to sit exactly on a
            % grid line. Anchoring on the spine side keeps that margin
            % exact on both pages of the spread; any leftover from
            % non-integer division falls on the outer edge instead,
            % where it isn't compared side-by-side against its mirror.
            \ifnum#5=0
                \draw[step=\gridspacing mm,black!\shadeval,line width=\strokewidth mm]
                    ({148.5mm-\marginright mm},{210mm-\vtopmargin mm}) grid
                    (\marginleft mm,\vbottommargin mm);
            \else
                \draw[step=\gridspacing mm,black!\shadeval,line width=\strokewidth mm]
                    (\marginleft mm,{210mm-\vtopmargin mm}) grid
                    ({148.5mm-\marginright mm},\vbottommargin mm);
            \fi
        \fi
        \ifnum#3=3% Dots
            % The dot nearest the spine sits exactly \dotspacing+\gutter
            % from the spine (mirrored on both pages); any leftover from
            % non-integer division falls on the outer edge instead. The
            % top dot sits exactly \dotspacing from the top border;
            % bottom uses \vbottommargin so dots don't collide with the
            % page number.
            \pgfmathtruncatemacro{\ndotcols}{floor((148.5-\dotmarginright-\dotmarginleft)/\dotspacing)}
            \pgfmathtruncatemacro{\ndotrows}{floor((210-\dotspacing-\vbottommargin)/\dotspacing)}
            \foreach \i in {0,...,\ndotcols} {
                \foreach \j in {0,...,\ndotrows} {
                    \ifnum#5=0
                        \def\dotx{{148.5mm-\dotmarginright mm-\i*\dotspacing mm}}
                    \else
                        \def\dotx{{\dotmarginleft mm+\i*\dotspacing mm}}
                    \fi
                    \fill[black!\shadeval]
                        (\dotx,{210mm-\dotspacing mm-\j*\dotspacing mm})
                        circle[radius=\dotsize mm];
                }
            }
        \fi

        % Page number - left or right based on position
        \ifdim#1<74mm % Left page
            \node[black!70] at (\marginleft mm,7mm) {\small #4};
        \else % Right page
            \node[black!70] at ({148.5mm-\marginright mm},7mm) {\small #4};
        \fi
    \end{scope}
}

% Binding/sewing guide on the spine between the two facing pages: a dashed
% line with end ticks, confined to the sewable zone (from \sewingmargin to
% 210mm-\sewingmargin) so it never touches the top/bottom edge.
% #1 = spine x-offset expression
\newcommand{\spinemarker}[1]{%
    \draw[black!45,densely dashed,line width=\strokewidth mm]
        ({#1},\sewingmargin mm) -- ({#1},{210mm-\sewingmargin mm});
    \draw[black!45,line width=\strokewidth mm]
        ({#1-2mm},\sewingmargin mm) -- ({#1+2mm},\sewingmargin mm);
    \draw[black!45,line width=\strokewidth mm]
        ({#1-2mm},{210mm-\sewingmargin mm}) -- ({#1+2mm},{210mm-\sewingmargin mm});
}

% Generate \numsheets A4 sheets of a single page style, each holding two
% consecutively-numbered A5 pages, with a \newpage between sheets. For
% single-sided (simplex) printing only - see \generatenotebookduplex for
% the double-sided version used to build an actual bound notebook.
% #1 = page type (1=ruled, 2=grid, 3=dots)
\newcommand{\generatenotebook}[1]{%
    \foreach \sheetnum in {1,...,\numsheets} {%
        \pgfmathtruncatemacro{\leftpagenum}{\startpage+2*\sheetnum-2}%
        \pgfmathtruncatemacro{\rightpagenum}{\leftpagenum+1}%
        \noindent\begin{tikzpicture}[remember picture]
            \path[use as bounding box] (0,0) rectangle (297mm,210mm);
            \makepage{\startx}{0mm}{#1}{\leftpagenum}{0}
            \makepage{\secondx}{0mm}{#1}{\rightpagenum}{1}
            \spinemarker{\secondx}
        \end{tikzpicture}%
        \ifnum\sheetnum<\numsheets\newpage\fi%
    }%
}

% Generate \numsheets A4 sheets (4 pages each: 2 front, 2 back) for
% double-sided (duplex) printing, using standard saddle-stitch booklet
% imposition so that after duplex printing, stacking sheet 1..\numsheets
% (sheet 1 on top, becoming the outermost/cover sheet), sewing through
% the stack along the centerline, and folding the whole stack in half,
% the pages read in correct sequential order.
%
% For sheet k (1 = outermost) of a \ntotal-page notebook:
%   front: [\startpage+\ntotal-2k+1 | \startpage+2k-2]
%   back:  [\startpage+2k-1         | \startpage+\ntotal-2k]
% (left | right, as printed - left is the spine-on-right page, right is
% the spine-on-left page, same as \generatenotebook).
%
% PDF pages come out in sheet order (sheet1-front, sheet1-back,
% sheet2-front, ...), which is the standard order a duplex printer
% expects: odd PDF pages print on sheet fronts, even on backs, of
% consecutive physical sheets.
%
% Whether your printer's duplex "long-edge"/"short-edge" binding setting
% produces correctly-aligned (not mirrored/upside-down) back pages for
% this landscape layout is printer/driver-specific and can't be predicted
% from here - print a short test run (e.g. \numsheets{2}) first and check
% the page numbers/orientation on the back before committing to the full
% print.
%
% The sewing-guide spine marker is only drawn on the back (inside) of
% each sheet - the front is the outward-facing cover/spread you don't
% see while sewing through the fold from inside.
% #1 = page type (1=ruled, 2=grid, 3=dots)
\newcommand{\generatenotebookduplex}[1]{%
    \pgfmathtruncatemacro{\ntotal}{4*\numsheets}%
    \foreach \sheetnum in {1,...,\numsheets} {%
        \pgfmathtruncatemacro{\frontleft}{\startpage+\ntotal-2*\sheetnum+1}%
        \pgfmathtruncatemacro{\frontright}{\startpage+2*\sheetnum-2}%
        \pgfmathtruncatemacro{\backleft}{\startpage+2*\sheetnum-1}%
        \pgfmathtruncatemacro{\backright}{\startpage+\ntotal-2*\sheetnum}%
        \noindent\begin{tikzpicture}[remember picture]
            \path[use as bounding box] (0,0) rectangle (297mm,210mm);
            \makepage{\startx}{0mm}{#1}{\frontleft}{0}
            \makepage{\secondx}{0mm}{#1}{\frontright}{1}
        \end{tikzpicture}%
        \newpage%
        \noindent\begin{tikzpicture}[remember picture]
            \path[use as bounding box] (0,0) rectangle (297mm,210mm);
            \makepage{\startx}{0mm}{#1}{\backleft}{0}
            \makepage{\secondx}{0mm}{#1}{\backright}{1}
            \spinemarker{\secondx}
        \end{tikzpicture}%
        \ifnum\sheetnum<\numsheets\newpage\fi%
    }%
}
